\subsection*{一步一步建立直观}

\paragraph{1. 把未知坐标看成两份面积贡献}

设输出 $b=xu+yv$，其中 $u,v$ 是矩阵两列。右图显示输出由这两根方向按系数组合得到。若只想求 $x$，就寻找一种测量，能让沿 $v$ 的那份贡献自动消失，而保留沿 $u$ 的贡献。这样的测量应该比较有向面积，而不是仅比较箭头长短。



\paragraph{2. 用替换列隔离一个系数}

取 $b$ 与 $v$ 的有向面积，沿 $v$ 的贡献因为共线而为零，于是面积为 $x$ 倍原底面面积。因此 $x=\det(b,v)/\det(u,v)$；换另一个底边同理得到 $y$。这个推导说明公式为何替换整列，而非单独替换某个元素。



\paragraph{3. 分母条件对应可见的几何障碍}

把两列调成共线后，分母为零。此时面积测量已经不能区分不同配方，继续做面积比没有意义。非零行列式不仅是为了避免除零，也是保证这套测量能够恢复两份坐标信息的实质条件。


